\section{引言}

开放飞行监视网络和研究档案正在持续改变航空数据分析的研究条件。OpenSky Network 证明了众包 ADS-B 传感网络可以长期支撑大规模可复现研究~\cite{opensky2014}，后续发布的 OpenSky 2019--2020 数据集又进一步说明，开放飞行数据已经能够覆盖大范围、长时间跨度的民航运行~\cite{strohmeier2021crowdsourced}。EUROCONTROL 的 Aviation Data Repository for Research（ADRR）则向研究和开发提供了更结构化的商业航班轨迹与运行信息档案~\cite{eurocontrol_adrr_2025}。在这样的数据条件下，航空轨迹预测不再只能围绕少量航班、单一机场或封闭基准展开，而可以被重新表述为一个面向开放数据、长时域和可复现实验的问题。

但数据规模的变化并不自动等于问题设定的变化。已有航空轨迹预测研究大多仍以连续坐标或四维轨迹回归为中心，例如 Attention-LSTM 四维轨迹预测模型~\cite{jia2022attentionlstm}、基于时频小波变换的轨迹预测方法~\cite{zhang2023wtftp}，以及更近期的直接多时域 FlightBERT++ 框架~\cite{guo_flightbertpp_2024}。另一些工作关注噪声 ADS-B 观测和稀疏卫星观测等更困难的数据条件~\cite{li2025noiserobust,shan2025sparse}。这些方法说明，更强的时间编码和多时域建模确实能够提升点式轨迹预测性能；但当研究对象变成开放运行档案上的 60--120 分钟预测时，单一路径坐标外推并不总是最合适的输出形式。

原因在于，开放运行数据上的长时域未来本身具有更强的不确定性和多模态性。相同的历史片段可能对应不同航路选择、不同阶段的爬升或下降过程，甚至不同的速度控制策略。此时，把模型能力集中在“给出一条最像真的未来轨迹”上，往往既不利于刻画不确定性，也不利于后续运行解释。对于战略交通态势感知、监视韧性分析、区域级流量推理和候选航路筛选而言，更自然的问题表述是：在给定预测提前量时，航空器未来最可能出现在哪些区域上，以及这些区域的概率是否可信。

因此，本文研究\emph{多时域区域航空轨迹预测}。与传统点预测不同，本文不是直接回归未来经纬度序列，而是在多个预测提前量上输出未来 H3 区域概率分布。这样的任务表述有三点实际意义。第一，它显式承认并表达长时域预测的不确定性，而不是把所有不确定性压缩进一条点轨迹。第二，它与开放轨迹档案的空间分辨率和运行多样性更匹配，避免把建模目标过度绑定到精确点位误差。第三，它天然支持面向概率质量和区域效率的评价，包括负对数似然、Brier 分数、期望校准误差，以及候选区域覆盖率、候选 cell 数量和候选面积。

本文工作的技术起点是一个已有的层级 H3 区域预测器，其核心为 GRU 历史编码器和 coarse-to-fine 解码头。我们保留这一层级空间输出框架，因为它已经能够把区域概率预测与父子层级约束、候选区域生成和可解释评估统一起来。本文进一步关注的是时间建模问题：在大规模开放航空轨迹上，什么样的时间结构更适合多时域概率区域预测？我们的实验表明，关键不在于简单地把编码器做得更深，而在于让不同预测时域能够从历史序列中读取不同的信息子集。

基于这一判断，本文提出\emph{多尺度 horizon-query Transformer}。模型由细粒度历史编码分支、粗粒度池化历史编码分支以及面向不同预测时域的可学习 query 组成；不同 horizon 通过 query attention 从编码状态中检索与本时域最相关的历史上下文，而后接入与基线兼容的层级 H3 解码器。这样的设计使时间建模改进可以在不破坏原有输出协议的前提下被单独评估，也便于与 Markov 基线、层级 GRU 基线、结构消融模型以及适配后的外部基线进行公平比较。

本文主要贡献如下：
\begin{enumerate}
  \item 基于 ADRR 派生的一个月商业航班轨迹数据，将大规模开放运行档案上的多时域航空预测问题形式化为层级 H3 概率区域预测任务，并给出可复现的数据划分与评估协议。
  \item 提出多尺度 horizon-query Transformer，在保持 coarse-to-fine H3 输出结构不变的前提下，显式建模不同预测时域对历史上下文的差异化需求，并在完整数据集上稳定提升区域准确率与概率质量。
  \item 建立兼顾精度、概率质量和区域效率的评估视角，同时报告 top-k 命中率、NLL、Brier 分数、ECE、候选覆盖率、候选 cell 数量与候选面积，而不是只依赖单一命中率指标。
  \item 给出与经典 Markov 基线、层级 GRU 基线、结构消融模型以及适配后的强外部基线的比较，并通过 H3 邻域 refinement 消融说明候选召回与概率校准之间存在明确权衡。
\end{enumerate}
